% Zitierstil-Konfiguration
% Der Stil wird in Diplomarbeit.tex mit \zitierstil ausgewählt:
%   htl        -> Zitierregeln für Diplomarbeiten (Amerikanische Zitierweise),
%                 im Text: (Hoffmann, 2020, S. 45)
%   alphabetic -> Kürzel aus dem shorthand-Feld, im Text: [PC:Web01]
% Erklärung siehe textparts/0/1_latex.tex, Abschnitt "Zitieren".

\usepackage{etoolbox}

\ifdefstring{\zitierstil}{htl}{%
  %---------------------------------------------------------------------------
  % HTL-Zitierregeln (Autor, Jahr)
  %---------------------------------------------------------------------------
  \usepackage[backend=biber, style=authoryear,
              maxcitenames=2, mincitenames=1, maxbibnames=99,
              uniquename=false, uniquelist=false, labeldate=year, mergedate=basic,
              date=short, urldate=short, isbn=false]{biblatex}
  % Zitate und Literaturverzeichnis auf Deutsch ("S.", "Hrsg.", Datum 02.12.2024),
  % auch wenn der Rest der Arbeit Englisch ist
  \DeclareLanguageMapping{english}{german}

  % shorthand und publisher werden ignoriert, damit dieselben .bib-Dateien
  % für beide Stile funktionieren (shorthand würde sonst (Autor, Jahr) ersetzen)
  \DeclareSourcemap{
    \maps[datatype=bibtex]{
      \map{
        \step[fieldset=shorthand, null]
        \step[fieldset=publisher, null]
      }
    }
  }

  % Im Text: (Hoffmann, 2020, S. 45)
  \renewcommand*{\nameyeardelim}{\addcomma\space}

  % Im Verzeichnis: Nachname, Vorname bei allen Autoren
  \DeclareNameAlias{sortname}{family-given}
  \DeclareNameAlias{default}{family-given}

  % Zwei Autoren mit "und", ab drei Autoren mit "; "
  \renewcommand*{\multinamedelim}{\addsemicolon\space}
  \renewcommand*{\finalnamedelim}{%
    \ifnumgreater{\value{liststop}}{2}{\addsemicolon\space}{\addspace und\space}}

  % Teile mit Komma trennen, Titel nicht kursiv, Aufsätze in Anführungszeichen
  \renewcommand*{\newunitpunct}{\addcomma\space}
  \DeclareFieldFormat*{title}{#1}
  \DeclareFieldFormat[article,incollection,inbook,inproceedings]{title}{\mkbibquote{#1}}
  \DeclareFieldFormat*{journaltitle}{#1}
  \DeclareFieldFormat*{booktitle}{#1}
  \DeclareFieldFormat*{maintitle}{#1}
  \DeclareFieldFormat{doi}{\url{https://doi.org/#1}}
  \DeclareFieldFormat{url}{\url{#1}}
  \DeclareFieldFormat{urldate}{\bibstring{urlseen}\space#1}
  \renewcommand*{\finentrypunct}{}               % kein Punkt am Ende

  % Zeitschrift: Name der Zeitschrift, 2/2023
  \renewbibmacro*{volume+number+eid}{%
    \setunit{\addcomma\space}%
    \printfield{volume}\setunit*{\adddot}\printfield{number}%
    \setunit{\bibeidpunct}\printfield{eid}}

  % Sammelband: In: Nachname, Vorname (Hrsg.), Titel; Zeitschrift ohne "In:"
  \renewbibmacro*{in:}{%
    \ifentrytype{article}{}{%
      \bibstring{in}\intitlepunct
      \ifnameundef{editor}{}{%
        \printnames{editor}\addspace\mkbibparens{\bibstring{editor}}%
        \clearname{editor}\newunit}}}

  % Internetquelle: Weblink, Typ, Zugriff am: Datum
  \renewbibmacro*{url+urldate}{%
    \usebibmacro{url}%
    \ifentrytype{online}{\newunit\printfield{addendum}\clearfield{addendum}}{}%
    \iffieldundef{urlyear}{}{\newunit\usebibmacro{urldate}}}

  % für englische (english -> german, siehe oben) und deutsche Arbeiten
  \newcommand{\htlbibstrings}[1]{%
    \DefineBibliographyStrings{#1}{
      edition   = {Auflage},
      volume    = {Band},
      in        = {In},
      urlseen   = {Zugriff am:},
      andothers = {et\addabbrvspace al\adddot},
    }}
  \forcsvlist{\htlbibstrings}{english,german,ngerman}
}{%
  %---------------------------------------------------------------------------
  % Kürzel (alphabetic), z.B. [PC:Web01]
  %---------------------------------------------------------------------------
  \usepackage[backend=biber, style=alphabetic]{biblatex}
}
